% DocInsights 2026 — Dr.DocBench (Track 2) system paper.
% Uses the official ACL template style files (acl.sty, acl_natbib.bst) — place this
% file alongside them (download from https://github.com/acl-org/acl-style-files).
\documentclass[11pt]{article}
\usepackage[final]{acl}

\usepackage{times}
\usepackage{latexsym}
\usepackage[T1]{fontenc}
\usepackage[utf8]{inputenc}
\usepackage{microtype}
\usepackage{inconsolata}
\usepackage{booktabs}
\usepackage{graphicx}
\usepackage{amsmath}
\usepackage{url}
\usepackage{enumitem}

\newcommand{\best}[1]{\textbf{#1}}

\title{Image Fidelity, Not Model Choice: A Simple Recipe for\\
Expert-Level Single-Page Document Parsing on Dr.DocBench}

\author{Author Name \\
  Affiliation \\
  \texttt{email@example.com}}

\begin{document}
\maketitle

\begin{abstract}
We describe our submission to Track~2 (Dr.DocBench) of the DocInsights~2026 shared
task, which converts single document-page images into structurally faithful
Markdown, scored on Text edit distance, Table TEDS, Formula CDM, and Reading Order.
Our system is a hosted vision--language model (VLM) pipeline: Qwen3.8-max as the
primary parser, with a Gemini~3.7-Flash ensemble routed to table pages. Rather than
switching models or engineering ever-larger prompts, we found that the single most
effective lever was \emph{image fidelity}: upscaling the low-resolution scans to the
VLM encoder's pixel budget and encoding them at maximum JPEG quality together
improved Overall by $+1.67$, twice the gain of every other change combined. On the
EvalAI Public Validation set (509 pages) our system reaches an Overall of
\best{71.33}. We report an extensive set of \emph{negative} results---reading-order
reordering, learned reading-order models, document-class conditioning, and
column-cropping all failed---and the lightweight validation discipline that let us
identify each failure on a held-out split before spending a leaderboard submission.
Code is publicly available.\footnote{\url{https://github.com/jinyichao/DocBench}}
\end{abstract}

\section{Introduction}
Expert-level document parsing---faithfully converting a scanned page into
structured Markdown---sits at the intersection of OCR, layout analysis, and
structured extraction. Dr.DocBench \citep{drdocbench2026}, the Track~2 benchmark of
the DocInsights~2026 shared task, is deliberately \emph{difficulty-aware}: its pages
were selected around cases where strong document parsers struggle, spanning dense
tables, irregular multi-column layouts, formulas, chemistry notation, and musical
scores. Each page is scored on four components---Text edit distance, Table TEDS,
Formula CDM, and Reading Order---converted to a $0$--$100$ scale and averaged per
page; the Overall is the mean of per-page scores over all scorable pages.

Our guiding question was pragmatic: given a strong hosted VLM, where does the
remaining Overall actually come from? We answer it with a disciplined ablation
campaign in which every candidate change was validated on a held-out development
split before consuming a leaderboard submission. The headline finding is that the
gains do \emph{not} come from model selection or prompt sophistication, both of
which we found largely saturated. They come from \textbf{feeding the model
higher-fidelity images}. The Dr.DocBench scans are low-resolution (median long edge
$968$\,px); at native resolution, dense text is under-tokenized by the VLM's
patch-based encoder. Upscaling each page to the encoder's pixel budget and encoding
at maximum JPEG quality raised Overall from $69.65$ to $71.33$.

Our contributions are: (i) a simple, reproducible pipeline that reaches Overall
$71.33$ on Dr.DocBench Public Validation; (ii) the empirical finding that image
fidelity---not model or prompt---is the dominant lever for low-resolution document
parsing, together with a measurement of the VLM's internal resolution ceiling;
and (iii) a documented set of negative results and the validation methodology that
produced them, which we believe is the most transferable outcome for future
participants.

\section{System Description}

\subsection{Engines}
The primary parser is \textbf{Qwen3.8-max} \citep{qwenvl} accessed through the
DashScope OpenAI-compatible endpoint, with internal ``thinking'' disabled (parsing
does not benefit from a hidden scratchpad, and disabling it is faster and cheaper).
We add a second engine, \textbf{Gemini~3.7-Flash} \citep{gemini}, used in two narrow
roles: (a) an ensemble
component routed to table pages, where we found it consistently stronger on TEDS,
and (b) a recovery path for pages that DashScope's account-level content filter
hard-blocks (returning \texttt{data\_inspection\_failed}), which we otherwise emit
as empty Markdown.

\subsection{Prompt Design}
Both engines share a single prompt, tuned to the scorer's output conventions rather
than to human readability. Key decisions, each validated on dev:
\begin{itemize}[nosep]
  \item \textbf{Tables} are emitted as raw HTML with \texttt{<td>} for every cell
    (never \texttt{<th>}, which the TEDS normaliser collapses anyway), the exact row
    $\times$ column grid, and no in-cell \texttt{<br>}; hyphen-split words are
    rejoined. This ``v2'' table block alone added ${\approx}5$ TEDS on dev over a
    naive HTML instruction.
  \item \textbf{Formulas}: isolated equations in display delimiters
    (\texttt{\$\$\ldots\$\$}); inline math in \texttt{\textbackslash(\ldots
    \textbackslash)}.
  \item \textbf{Chemistry}: \texttt{\textbackslash ce\{\ldots\}} for expressions and
    fenced \texttt{smiles} blocks for drawn 2D structures. The scorer treats
    chemistry as \emph{text}, so we do not wrap it in math delimiters.
  \item \textbf{Reading order}: column-major (finish the left column before the
    right). We verified on dev that ground-truth reading order is column-major on
    $35$ of $38$ clean two-column pages, matching the model's default.
\end{itemize}

\subsection{Image Encoding: The Core Lever}
\label{sec:image}
Our encoder resizes each page so the long edge is $2000$\,px (upscaling low-resolution
scans, not only downscaling large ones) and encodes to JPEG at quality~$100$ with
chroma subsampling disabled. Two mechanisms explain the gains:

\paragraph{Clearing the visual-token floor.} VLM encoders tile an image into
fixed-size patches; text that spans too few patches cannot be resolved. Upscaling a
$968$\,px scan to $2000$\,px raises the number of visual tokens devoted to the same
(interpolated) glyphs, letting the encoder separate characters it previously smeared
together. This adds no new pixel information but removes a tokenisation bottleneck.

\paragraph{The resolution ceiling.} Qwen applies an internal \texttt{max\_pixels}
cap. We measured the billed input tokens as a function of the long edge
(Table~\ref{tab:tokens}): tokens grow up to $2000$\,px and then plateau---$2500$\,px
and $3000$\,px yield the \emph{identical} token count, i.e.\ the model silently
downscales them to the same image. Thus $2000$\,px is exactly the sweet spot, and
input resolution beyond it is discarded. Requesting a larger budget server-side
(\texttt{vl\_high\_resolution\_images}) does raise the token count but did
\emph{not} improve scores, because the native scans carry no detail beyond the cap.
Within the fixed pixel budget, however, JPEG quality still matters: raising quality
from $90$ to $100$ reduces compression artefacts on small glyphs and improved both
Text and Reading Order on the test set.

\begin{table}[t]
\centering
\small
\begin{tabular}{lccccc}
\toprule
Long edge (px) & 1400 & 1600 & 2000 & 2500 & 3000 \\
\midrule
Input tokens & 2323 & 2771 & \best{3423} & 3423 & 3423 \\
\bottomrule
\end{tabular}
\caption{Billed Qwen input tokens per page vs.\ input resolution (mean over dense
dev pages). The token count plateaus at $2000$\,px: higher input resolution is
silently downscaled to the same image.}
\label{tab:tokens}
\end{table}

\subsection{Ensemble Routing and Assembly}
A layout detector \citep[DocLayout-YOLO;][]{doclayout} flags table pages; on those pages we replace the
Qwen output with Gemini's when Gemini emits a valid, non-truncated \texttt{<table>}
of comparable length (a guard against under-generation). Content-filter--blocked
pages are filled from the Gemini recovery run. The remaining pages keep the Qwen
output. The assembled tree is packaged into the required \texttt{predictions.jsonl}
after validation against the manifest.

\subsection{Robustness}
Larger upscaled payloads make hosted calls prone to hanging connections; we set a
per-request timeout so a stalled call aborts and retries, and wrap batch inference in
a self-healing supervisor that restarts on lack of progress until every page is
complete. Inference is resume-safe (a page is skipped if its Markdown already
exists).

\section{Experiments}

\subsection{Setup}
We evaluate on the EvalAI Public Validation phase ($509$ required pages, of which
$487$ are scorable). For local iteration we use a held-out development split with
released ground truth, reusing the official DrDocBench metrics (Text edit distance,
TEDS \citep{teds}, reading-order edit distance) and an offline Formula CDM
\citep{cdm} renderer
(\LaTeX\,$\rightarrow$\,PDF\,$\rightarrow$\,Hungarian character matching). All
design decisions were made against dev; the leaderboard was used only to confirm
banked improvements.

\subsection{Main Result and Ablation}
Table~\ref{tab:ablation} traces the Overall on the leaderboard as each validated
lever was added. The two image-fidelity steps (upscaling, JPEG~q100) account for
$+1.67$ of the total $+2.07$ improvement over the initial submission.

\begin{table}[t]
\centering
\small
\begin{tabular}{lc}
\toprule
System & Overall \\
\midrule
Qwen + v2 table prompt & 69.26 \\
\;+ Gemini table ensemble & 69.65 \\
\;+ image upscaling to 2000\,px & 70.68 \\
\;+ \texttt{smiles} for drawn structures & 70.69 \\
\;+ JPEG quality 100 on low-res pages & \best{71.33} \\
\bottomrule
\end{tabular}
\caption{Leaderboard ablation. Each row adds one dev-validated change. Image fidelity
(rows 3 and 5) is the dominant contributor.}
\label{tab:ablation}
\end{table}

\subsection{Component Comparison}
Table~\ref{tab:compare} places our per-component dev scores alongside the reference
baselines reported by the benchmark. On shared components our profile is competitive
to strong: best Text and Formula-Edit, highest Formula CDM, and a Table TEDS that
matches the best specialised parser. We stress that the \emph{Overall} is not
directly comparable: our dev split is prose-heavy and, as we discuss below, contains
essentially no isolated display formulas, so its composition differs from the
reference benchmark. The component cells, not the Overall, are the meaningful
comparison.

\begin{table*}[t]
\centering
\small
\begin{tabular}{lcccccc}
\toprule
Model & Text\,$\downarrow$ & Formula Edit\,$\downarrow$ & CDM\,$\uparrow$ & TEDS\,$\uparrow$ & Order\,$\downarrow$ & Overall\,$\uparrow$ \\
\midrule
MinerU 2.5 & 0.33 & 0.51 & 24.15 & 55.85 & 0.30 & 54.37 \\
PaddleOCR & 0.73 & 0.42 & 30.73 & 51.79 & 0.71 & 34.78 \\
Kimi-K2.5 & 0.19 & 0.35 & 27.04 & 51.98 & 0.18 & 60.38 \\
Claude Opus 4.6 & 0.22 & 0.37 & 32.02 & 49.21 & 0.19 & 60.19 \\
Gemini 3.1 Pro & 0.22 & 0.35 & 32.22 & 51.26 & 0.21 & 60.13 \\
GPT-5.5 & 0.19 & 0.36 & 34.55 & 48.90 & 0.17 & 61.94 \\
\midrule
Ours (dev) & \best{0.11} & \best{0.19} & \best{55.39} & \best{55.89} & \best{0.12} & 86.81$^{\dagger}$ \\
\bottomrule
\end{tabular}
\caption{Per-component comparison against reported baselines. Best per column in
\textbf{bold}. $^{\dagger}$Our Overall is on a prose-heavy dev split with a different
composition (notably, no display formulas) and is \emph{not} a head-to-head number;
the component cells are the meaningful comparison.}
\label{tab:compare}
\end{table*}

\section{Analysis and Lessons Learned}

\paragraph{Image fidelity dominates.} For low-resolution scans, the bottleneck is
how much of the page the encoder can resolve, not the model's parsing ability. The
two fidelity steps outweighed every prompt and ensemble change combined, and the
resolution ceiling (Table~\ref{tab:tokens}) tells us precisely when the lever is
exhausted.

\paragraph{Exploiting scorer structure.} The Text metric penalises \emph{missing}
ground-truth content but not \emph{extra} predicted blocks (unmatched predictions
are discarded, not penalised). This asymmetry justifies an ``emit when unsure''
policy---e.g.\ transcribing text embedded in figures, or emitting a \texttt{smiles}
block for every drawn structure---since surplus content is free while omission is
costly.

\paragraph{Negative results.} Several intuitive levers failed, each caught on dev:
\begin{itemize}[nosep]
  \item \textbf{Reading order} resisted every approach we tried: a stronger
    column-discipline prompt (net-negative), geometric XY-cut ($78.4$ even on
    ground-truth boxes, below the model's $81.85$), an off-the-shelf learned
    reading-order model \citep[LayoutReader;][]{layoutreader} (far below both), and routing to Gemini
    (worse on non-table multi-column prose). Ground-truth order depends on page
    \emph{content}, which the VLM already sees; a box-only model cannot beat it.
  \item \textbf{Formula CDM} is near-zero on this test set---but a layout-detector
    sweep confirmed the Public Validation split contains essentially \emph{no}
    isolated display formulas, so the metric is a near-empty sample, not a failure.
    Forcing display-math would only mis-categorise inline chemistry.
  \item \textbf{Document-class conditioning} and same-class few-shot exemplars
    appeared to help but were traced to same-document leakage; the honest
    cross-document effect was neutral.
  \item \textbf{Column cropping} (detect columns, transcribe each separately)
    regressed, because detector false positives split single-column pages.
\end{itemize}

\paragraph{Validation discipline.} The recurring lesson is that a cheap held-out
evaluation---matched to the deployment condition (detector-selected, not
oracle-selected, pages; cross-document, not same-document, exemplars)---is worth far
more than leaderboard probing. Every negative above was identified on dev before
spending a submission, and two would have been actively harmful had we deployed them
on the basis of a flawed validation.

\section{Conclusion}
Our Dr.DocBench system reaches Overall $71.33$ with a deliberately simple recipe: a
strong hosted VLM, a scorer-tuned prompt, a narrow table-page ensemble, and---above
all---high-fidelity image encoding calibrated to the model's resolution ceiling. The
broader message for expert-level document parsing on low-resolution scans is that
the image, not the model, is often the binding constraint. We release our code and a
record of the negative results to help future participants spend their effort where
it pays off.

\section*{Limitations}
Our Overall is reported on the Public Validation split, whose composition (prose-heavy,
no display formulas) differs from the full benchmark, so component scores are more
informative than the aggregate. The pipeline depends on two proprietary hosted APIs,
whose behaviour (content filtering, internal \texttt{max\_pixels} cap) may change and
is not fully controllable. Reading order remains our weakest scorable component, and
we did not pursue model fine-tuning.

\bibliography{references}

\end{document}
